\section*{Capítulo \ref{ch:b3:plant-molecular-physiology} --- Fisiología molecular vegetal y respuestas al estrés}

\begin{solution}{exo:b3:plant-molecular-physiology:1}
\omterm{def:b3:plant-molecular-physiology:photoreceptors}{Fitocromos}, rojo \qty{660}{nm} y rojo lejano \qty{730}{nm}:
desetiolación, \omterm{prop:b3:plant-molecular-physiology:photoequilibrium}{evitación de la sombra}, germinación. \omterm{def:b3:plant-molecular-physiology:photoreceptors}{Criptocromos}, azul
\qty{450}{nm}: desetiolación, reloj, floración. \omterm{def:b3:plant-molecular-physiology:photoreceptors}{Fototropinas}, azul:
curvatura hacia la luz, apertura estomática, movimiento de los
cloroplastos. La clorofila mide la luz como energía y no hace nada con
la información; los fotorreceptores leen el color, la dirección y la
duración a flujos un millón de veces menores y cambian la expresión
génica ---la diferencia entre comerse la luz y leerla.
\end{solution}

\begin{solution}{exo:b3:plant-molecular-physiology:2}
Cada \omterm{def:b3:endocrinology:hormone}{hormona} promueve la ubiquitinación y la destrucción de un
represor: la \omterm{def:b3:plant-molecular-physiology:hormones}{auxina} pega \omterm{def:b3:plant-molecular-physiology:hormones}{TIR1} a las proteínas Aux/IAA, la \omterm{def:b3:plant-molecular-physiology:hormones}{giberelina}
unida a GID1 entrega las proteínas \omterm{def:b3:plant-molecular-physiology:hormones}{DELLA} a una ligasa F-box, y el
jasmonato pega COI1 a las proteínas JAZ. Los represores ya están
sentados sobre sus genes diana con los activadores preparados al lado,
de modo que no hace falta fabricar proteína nueva: destruir el
represor, cosa que el proteasoma hace en minutos, enciende los genes de
inmediato.
\end{solution}

\begin{solution}{exo:b3:plant-molecular-physiology:3}
El ABA se une a PYR/PYL; el complejo inhibe la fosfatasa PP2C; la
cinasa SnRK2 (OST1), ya no desfosforilada, se activa y fosforila el
canal aniónico SLAC1; los aniones salen, la membrana se despolariza,
se abren los canales de salida de potasio, salen el K$^{+}$ y después
el agua, la turgencia cae y las \omterm{prop:b3:plant-molecular-physiology:stomata}{células oclusivas} se desinflan juntas.
Cualquier eslabón puede romperse: sin receptores, con una fosfatasa que
no puede inhibirse (el clásico mutante \emph{abi1}), sin OST1 o sin
SLAC1 ---cada caso da una planta lacia cuyos estomas siguen abiertos en
la sequía.
\end{solution}

\begin{solution}{exo:b3:plant-molecular-physiology:4}
La PTI reconoce moléculas comunes a clases microbianas enteras
(flagelina, quitina) con cinasas receptoras de la superficie celular y
termina en una pared reforzada, estomas cerrados, expresión de genes de
defensa y un crecimiento más lento, por lo general sin muerte celular.
La ETI reconoce una proteína efectora concreta o su daño con un
\omterm{def:b3:plant-molecular-physiology:immunity}{receptor NLR} intracelular y termina en la muerte hipersensible de las
células infectadas y en una alarma sistémica.
\end{solution}

\begin{solution}{exo:b3:plant-molecular-physiology:5}
Luz del día, $0.87\times 1.15/1.75 = 0.57$; una hoja, $0.87\times 0.4/1.0 =
0.35$; cubierta densa, $0.87\times 0.1/0.7 = 0.12$; lámpara
incandescente, $0.87\times 0.7/1.3 = 0.47$. Una lámpara de filamento es
rica en rojo lejano, de modo que $\varphi$ queda por debajo del de la
luz del día y la plántula lee sombra parcial: se alarga. Las lámparas
blancas frías y los LED sin rojo lejano hacen lo contrario.
\end{solution}

\begin{solution}{exo:b3:plant-molecular-physiology:6}
Pared $5.0$, citosol $7.5$: $(1 + 10^{2.75})/(1 + 10^{0.25}) = 563/2.78
= 203$. Pared $4.5$: $563/(1 + 10^{-0.25}) = 563/1.56 = 360$
---acidificar la pared eleva la fracción neutra de fuera y con ella la
captación. Citosol $6.5$ con pared $5.5$:
$(1 + 10^{1.75})/6.62 = 57/6.62 = 8.6$: el atrapamiento cae cinco
veces, la \omterm{def:b3:plant-molecular-physiology:hormones}{auxina} se escapa por todas las caras y los gradientes
polares se aplanan ---una razón de que las raíces en anoxia pierdan la
orientación.
\end{solution}

\begin{solution}{exo:b3:plant-molecular-physiology:7}
\qty{1}{cm/h} $= \qty{10000}{\micro m/h}$: $200$ células por hora,
\qty{18}{s} en cada una. Difusión a través de \qty{50}{\micro m}: $t =
(5\times 10^{-5})^{2}/(2\times 5\times 10^{-10}) = \qty{2.5}{s}$. El
interior de la célula se cruza en segundos; la mayor parte de los
\qty{18}{s} se pasa esperando a ser sacada por los PIN de la membrana
basal, que es el paso limitante.
\end{solution}

\begin{solution}{exo:b3:plant-molecular-physiology:8}
$E/A = 1.6\,\Delta w/\Delta c = 1.6\times\num{16000}/150 = 171$
moléculas de agua por CO$_{2}$. C$_{4}$, con $\Delta c = 250$: $102$.
Día seco, con $\Delta w = 24$: $256$.
\end{solution}

\begin{solution}{exo:b3:plant-molecular-physiology:9}
La congelación deshidrata las células al llevar el agua al hielo
extracelular, de modo que el frío y la sequía son ambos estreses de
deshidratación. Comparten segundos mensajeros (ABA, picos de calcio,
oxígeno reactivo), factores de transcripción (la familia CBF/DREB la
inducen los dos) y productos: dehidrinas y proteínas LEA que protegen
membranas y proteínas de la pérdida de agua, solutos compatibles
(prolina, azúcares) que retienen agua, y antioxidantes. Una planta que
los ha fabricado para un estrés los tiene para el otro.
\end{solution}

\begin{solution}{exo:b3:plant-molecular-physiology:10}
Con el \qty{1}{\%} del sol pleno, $k_{1} + k_{2} = \qty{0.005}{s^{-1}}$:
el 95~\% del equilibrio tras $3/k = \qty{600}{s}$, diez minutos, frente
a seis segundos al mediodía. Ocho horas son cuatro semividas: $1/16 \approx
\qty{6}{\%}$ del Pfr queda. Un destello breve fija $\varphi$, pero el
Pfr decae después, mientras que un día lo mantiene alto durante horas;
las respuestas que exigen que el Pfr actúe durante horas (germinación,
desetiolación) integran la exposición, de modo que una semilla sacada
un instante por un arado no responde como una que yace en la
superficie.
\end{solution}

\begin{solution}{exo:b3:plant-molecular-physiology:11}
Diez focos al día a $100$ células cada uno: \num{1000} células,
$10^{-4}$ de la hoja; a lo largo de una temporada, una fracción de un
uno por ciento, frente a la pérdida de la hoja si una sola infección se
extendiera sin freno. Las células muertas no cuestan casi nada y se
llevan consigo el alimento del patógeno ---para un biótrofo, que
necesita células vivas. Un necrótrofo se alimenta de tejido muerto: la
\omterm{def:b3:plant-molecular-physiology:immunity}{respuesta hipersensible} le entrega una comida y un punto de apoyo, y
algunos (\emph{Botrytis}, \emph{Sclerotinia}) segregan toxinas que la
provocan a propósito; contra ellos la planta se apoya en las defensas
mediadas por el jasmonato y no en la muerte celular.
\end{solution}

\begin{solution}{exo:b3:plant-molecular-physiology:12}
Producción a un ritmo $p$ de la hora 12 a la 16. Día largo, con luz
todo el rato, $k = \ln 2/4 = \qty{0.17}{h^{-1}}$: el nivel en la 16 es
$(p/k)(1 -
\mathrm{e}^{-4k}) = (p/0.17)(0.5) = 2.9p$. Día corto, con oscuridad
desde la hora 10, $k = \qty{1.39}{h^{-1}}$: $(p/1.39)(1 - \mathrm{e}^{-5.5})
\approx 0.72p$. Cuatro veces más CONSTANS en el día largo. Un umbral
intermedio ---pongamos $2p$, necesario para activar \emph{FT}--- solo
se cruza cuando la ventana de producción del reloj coincide con la luz:
la coincidencia de un ritmo interno con el día externo convierte una
duración graduada del día en una decisión de florecer de todo o nada.
\end{solution}

\begin{solution}{pb:b3:plant-molecular-physiology:1}
\textbf{1.} Luz del día, $0.87\times 1.15/1.75 = 0.57$; cubierta, $0.87\times
0.2/0.8 = 0.22$.
\textbf{2.} $r = 2(1 - \varphi)$: \qty{0.86}{mm/h} a la luz del día y
\qty{1.57}{mm/h} a la sombra; en \qty{48}{h}: $41$ frente a
\qty{75}{mm}, \qty{34}{mm} de más.
\textbf{3.} El rojo cae a $0.1\times 1.15 = 0.115$ del rojo lejano, que
cae a $0.8$: $\zeta = 0.144$, $\varphi = 0.87\times 0.144/0.744 =
0.17$. Sí: una sola hoja baja $\varphi$ de $0.57$ a $0.17$, de lleno en
el intervalo de \omterm{prop:b3:plant-molecular-physiology:photoequilibrium}{evitación de la sombra}.
\textbf{4.} Las dos velocidades de fotoconversión escalan con el flujo,
de modo que su razón no. Bajo las nubes el espectro apenas cambia y
$\varphi$ se queda cerca de $0.57$: una plántula a cielo abierto en un
día gris no se alarga ---lee vecinos, no brillo.
\textbf{5.} El \qty{95}{\%} tras $3/(k_{1} + k_{2})$: \qty{6}{s} a pleno
sol y \qty{600}{s} con el \qty{1}{\%}.
\textbf{6.} \qty{8}{h} $=$ cuatro semividas, $1/16 = \qty{6.3}{\%}$;
\qty{14}{h} $=$ siete, $1/128 = \qty{0.8}{\%}$. El Pfr residual al
amanecer podría informar de la duración de la noche, pero la reversión
es una reacción térmica que corre más deprisa en las noches cálidas, y
el valor de partida depende de la luz del día: las plantas miden la
duración de la noche con el reloj.
\textbf{7.} Pared: $1/(1 + 10^{0.75}) = 0.15$; citosol: $1/(1 +
10^{2.45}) = 0.0035$.
\textbf{8.} $(1 + 282)/(1 + 5.6) = 283/6.6 = 43$. Citosol,
$43\times 0.1 = \qty{4.3}{\micro M}$.
\textbf{9.} $\qty{1}{cm/h}/\qty{100}{\micro m} = 100$ células por hora,
\qty{36}{s} en cada una.
\textbf{10.} \qty{5}{h}. La curvatura empieza en menos de una hora
porque la redistribución que importa es lateral, a lo ancho de un
milímetro de ápice, y la respuesta es local; la corriente de larga
distancia fija el suministro, no el momento.
\textbf{11.} Flanco inferior: $60/50 = 1.2$, \omterm{def:b3:plant-molecular-physiology:hormones}{auxina} un \qty{20}{\%} por
encima de la norma, elongación $0.8$; flanco superior: $0.8$ de la
norma, elongación $1.2$. El flanco superior crece más que el inferior:
la raíz se curva hacia abajo.
\textbf{12.} Las células del brote están por debajo de su óptimo de
\omterm{def:b3:plant-molecular-physiology:hormones}{auxina}, de modo que la \omterm{def:b3:plant-molecular-physiology:hormones}{auxina} de más del flanco inferior favorece su
crecimiento: el flanco inferior crece más que el superior y el brote se
curva hacia arriba.
\textbf{13.} $E = 0.3\times 16 = \qty{4.8}{mmol\,m^{-2}\,s^{-1}}$.
\textbf{14.} $A = (0.3/1.6)\times 150 = \qty{28}{\micro mol\,m^{-2}\,s^{-1}}$;
$E/A = 4800/28 = 171$.
\textbf{15.} Superficie \qty{0.002}{m^2}, \num{43200}~s: agua, $4.8\times
10^{-3}\times 0.002\times\num{43200} = \qty{0.41}{mol} = \qty{7.5}{g}$;
CO$_{2}$, $28\times 10^{-6}\times 0.002\times\num{43200} =
\qty{2.4e-3}{mol}$, es decir, $2.4\times 10^{-3}/6\times 180 =
\qty{0.073}{g}$ de glucosa ---cien gramos de agua por un gramo de
azúcar.
\textbf{16.} $E = 0.05\times 16 = \qty{0.8}{mmol\,m^{-2}\,s^{-1}}$,
$A = \qty{4.7}{\micro mol\,m^{-2}\,s^{-1}}$, razón todavía $171$. La
planta ha ahorrado seis séptimos de su agua y ha renunciado a seis
séptimos de su carbono: cerrar la válvula cambia el ritmo, no el
precio.
\textbf{17.} C$_{4}$: $\Delta c = 250$, $E/A = 4800/46.9 = 102$. CAM de
noche: $E = 0.3\times 5 = \qty{1.5}{mmol}$, $A = 28$: $E/A = 53$.
\textbf{18.} $E/A = 1.6\,\Delta w/\Delta c$: la conductancia se cancela
porque los dos gases pasan por el mismo poro. La planta puede cambiar
los gradientes ---concentrar el CO$_{2}$ por dentro (C$_{4}$), abrir
cuando el aire es húmedo (CAM, el amanecer), enfriar la hoja, engrosar
la capa límite con pelos---, pero no la física de dos gases que
comparten un agujero.
\textbf{19.} \num{1000} células al día, $10^{-4}$ de la hoja.
\textbf{20.} \num{60000} células en 60 días, un \qty{0.6}{\%}. Un
escape diario que destruyera el \qty{5}{\%} consumiría la hoja en
veinte días: la \omterm{def:b3:plant-molecular-physiology:immunity}{respuesta hipersensible} cuesta la centésima parte de lo
que cuesta una sola infección escapada.
\textbf{21.} $0.2\times 5 = \qty{1}{\%}$ de la fotosíntesis. Mantenida
siempre encendida costaría un \qty{5}{\%} y la planta quedaría
adelantada por vecinas que lo gastan en hojas; una defensa inducible
solo se paga cuando la amenaza está presente.
\textbf{22.} Gana invisibilidad ante ese NLR e infecta la variedad
resistente; pierde lo que hiciera el efector (suprimir la PTI), de modo
que es más débil en las plantas sin ese NLR. La población de plantas
pasa a favorecer NLR contra los efectores restantes, y el viejo gen de
resistencia, ya inútil, decae: un ciclo dependiente de la frecuencia en
los genes de ambos bandos.
\textbf{23.} El imitador (la coronatina) activa la vía del jasmonato,
que antagoniza la vía del salicilato; el salicilato es la defensa
contra los biótrofos, como la propia bacteria, y el imitador vuelve
además a abrir los estomas que la planta había cerrado. El patógeno
vuelve un brazo del sistema inmunitario contra el otro.
\textbf{24.} Un solo gen de resistencia repartido por millones de
plantas idénticas es una selección uniforme: cualquier mutante que
pierda o altere el efector se extiende por todo el cultivo en unas
pocas campañas. Los fitomejoradores apilan varios genes de resistencia
para que haya que perder varios efectores a la vez, mezclan variedades,
rotan los genes a lo largo de los años y favorecen la resistencia
amplia desencadenada por patrones, a la que ninguna mutación única
escapa.
\textbf{25.} $\varphi \approx 0.22$ bajo la cubierta; \omterm{def:b3:plant-molecular-physiology:hormones}{auxina} dentro :
fuera $\approx 43$; unas $170$ moléculas de agua perdidas por cada
CO$_{2}$ fijado.
\end{solution}
